\documentclass{article}
\usepackage{amsmath,amssymb,amsthm}
\usepackage{algorithm}
\usepackage{algpseudocode}
\usepackage{hyperref}
\usepackage{enumitem}
\newtheorem{theorem}{Theorem}
\newtheorem{lemma}{Lemma}
\newtheorem{assumption}{Assumption}
\newcommand{\E}{\mathbb{E}}
\newcommand{\norm}[1]{\left\|#1\right\|}
\newcommand{\ip}[2]{\langle #1,#2\rangle}
\begin{document}

\section*{Error‑Feedback Stochastic Gradient Descent (EF‑SGD)}

\section*{Mathematical Formulation}
Let $f:\mathbb{R}^d\to\mathbb{R}$ be the loss function.  
At iteration $k$ we have a current iterate $w_k\in\mathbb{R}^d$ and an
error buffer $e_k\in\mathbb{R}^d$ (initially $e_0=0$).  
A stochastic gradient $g_k$ is obtained:
\[
g_k \;=\; \nabla f(w_k;\xi_k),
\]
where $\xi_k$ is a random data sample.  
Define the \emph{pre‑quantized} vector
\[
u_k \;:=\; g_k+e_k .
\]
A (possibly biased) quantizer $Q:\mathbb{R}^d\to\mathbb{R}^d$ is applied:
\[
\tilde g_k \;:=\; Q(u_k).
\]
The error buffer is updated by the residual:
\[
e_{k+1}\;:=\; u_k-\tilde g_k = g_k+e_k-Q(g_k+e_k). \tag{1}
\]
Finally the model is updated with learning rate $\eta>0$:
\[
w_{k+1}\;:=\; w_k-\eta \tilde g_k . \tag{2}
\]

\section*{Pseudocode}
\begin{algorithm}[H]
\caption{EF‑SGD}
\begin{algorithmic}[1]
\State \textbf{Input:} initial point $w_0$, step size $\eta$, quantizer $Q$
\State $e_0\gets 0$
\For{$k=0,1,\dots,T-1$}
    \State Sample $\xi_k$ and compute stochastic gradient $g_k=\nabla f(w_k;\xi_k)$
    \State $u_k \gets g_k+e_k$
    \State $\tilde g_k \gets Q(u_k)$
    \State $e_{k+1}\gets u_k-\tilde g_k$ \Comment{Residual error (1)}
    \State $w_{k+1}\gets w_k-\eta \tilde g_k$ \Comment{Model update (2)}
\EndFor
\State \textbf{Return} $w_T$
\end{algorithmic}
\end{algorithm}

\section*{Dry Run Example}
Consider the scalar quadratic $f(w)=(w-3)^2$.
Its exact gradient is $\nabla f(w)=2(w-3)$.  
We use a deterministic 1‑bit quantizer that rounds to the nearest
multiple of $0.1$:
\[
Q(x)=0.1\cdot\operatorname{round}\!\left(\frac{x}{0.1}\right).
\]
Set $w_0=0$, $e_0=0$, and $\eta=0.1$.

\begin{center}
\begin{tabular}{c|c|c|c|c}
$k$ & $g_k=2(w_k-3)$ & $u_k=g_k+e_k$ & $\tilde g_k=Q(u_k)$ & $w_{k+1}=w_k-\eta\tilde g_k$ \\ \hline
0 & $-6$ & $-6$ & $-6.0$ & $0-0.1(-6)=0.6$ \\
1 & $2(0.6-3)=-4.8$ & $-4.8+0$ & $-4.8$ & $0.6-0.1(-4.8)=1.08$ \\
2 & $2(1.08-3)=-3.84$ & $-3.84+0$ & $-3.8$ & $1.08-0.1(-3.8)=1.46$
\end{tabular}
\end{center}
Objective values: $f(0)=9$, $f(0.6)=5.76$, $f(1.08)=3.6864$, $f(1.46)=2.3716$.
The error buffer stays zero because the quantizer is exact on the
displayed multiples.

\newpage

\section*{Assumptions}
\begin{assumption}[Smoothness]\label{ass:smooth}
$f$ is $L$‑smooth:
\[
\forall x,y\in\mathbb{R}^d,\qquad
f(y)\le f(x)+\ip{\nabla f(x)}{y-x}+\frac{L}{2}\norm{y-x}^2 .
\]
\end{assumption}
\begin{assumption}[Unbiased Stochastic Gradient]\label{ass:unbiased}
For all $k$, $\E_{\xi_k}[g_k\mid w_k]=\nabla f(w_k)$ and
\[
\E_{\xi_k}\!\big[\norm{g_k-\nabla f(w_k)}^2\mid w_k\big]\le\sigma^2 .
\]
\end{assumption}
\begin{assumption}[Quantizer Contraction]\label{ass:quant}
There exists $\delta\in(0,1]$ such that for any $x\in\mathbb{R}^d$,
\[
\E_Q\!\big[\norm{Q(x)-x}^2\big]\le (1-\delta)\norm{x}^2 .
\]
Moreover $Q$ is unbiased: $\E_Q[Q(x)]=x$.
\end{assumption}
\begin{assumption}[Bounded Second Moment]\label{ass:bound}
There exists $G>0$ with $\E\!\big[\norm{g_k}^2\big]\le G^2$ for all $k$.
\end{assumption}

\section*{Main Convergence Theorem}
\begin{theorem}\label{thm:main}
Let Assumptions~\ref{ass:smooth}–\ref{ass:bound} hold.
Choose a constant stepsize 
\[
\eta\le\frac{1}{L(1+\tfrac{1-\delta}{\delta})}
\;=\;\frac{\delta}{L(2-\delta)} .
\]
Denote $w^{*}$ as a (possibly non‑unique) stationary point.
Then for any $T\ge1$ the iterates produced by EF‑SGD satisfy
\[
\frac{1}{T}\sum_{k=0}^{T-1}\E\!\big[\norm{\nabla f(w_k)}^2\big]
\;\le\;
\underbrace{\frac{2\big(f(w_0)-f(w^{*})\big)}{\eta T}}_{\text{optimization term}}
\;+\;
\underbrace{\frac{L\eta G^2}{\delta}}_{\text{variance+quantization term}} .
\tag{3}
\]
Consequently, with $\eta = \Theta\!\big(\tfrac{1}{\sqrt{T}}\big)$,
\[
\min_{0\le k<T}\E\!\big[\norm{\nabla f(w_k)}^2\big]=O\!\left(\frac{1}{\sqrt{T}}\right).
\]
\end{theorem}

\section*{Theoretical Properties}
\begin{itemize}[leftmargin=*]
\item \textbf{Stability.} The error feedback recursion (1) guarantees that
the accumulated quantization error does not diverge; Lemma~\ref{lem:error}
shows $\norm{e_k}^2$ is uniformly bounded.
\item \textbf{Hyper‑parameter Sensitivity.}
The contraction factor $\delta$ appears only in the denominator of the
second term of (3); a larger $\delta$ (more accurate quantizer) yields
a tighter bound.
\item \textbf{Comparison to Baselines.}
Without error feedback ($e_k\equiv0$) the bound would contain the factor
$(1-\delta)^{-1}$, leading to a worse dependence on the quantizer quality.
EF‑SGD recovers the standard SGD rate $O(1/\sqrt{T})$ when $\delta=1$
(i.e., perfect communication).
\end{itemize}

\section*{Discussion}
The theorem shows that EF‑SGD attains the same asymptotic
convergence rate as vanilla SGD for non‑convex smooth objectives,
while only incurring an additive constant proportional to $1/\delta$.
Thus, even coarse quantization (small $\delta$) does not destroy the
$O(1/\sqrt{T})$ decay, provided the stepsize is chosen according to the
stability condition in the theorem.

\newpage

\section*{Preliminaries (Definitions and Lemmas)}
\begin{lemma}[Smoothness Inequality]\label{lem:smooth}
If $f$ is $L$‑smooth, then for any $x,y$,
\[
f(y)\le f(x)+\ip{\nabla f(x)}{y-x}+\frac{L}{2}\norm{y-x}^2 .
\]
\end{lemma}
\begin{lemma}[Error‑Feedback Recursion]\label{lem:error}
Under Assumption~\ref{ass:quant},
\[
\E\!\big[\norm{e_{k+1}}^2\mid\mathcal{F}_k\big]
\le (1-\delta)\E\!\big[\norm{e_k+\nabla f(w_k)}^2\mid\mathcal{F}_k\big]
    + (1-\delta)\sigma^2 ,
\]
where $\mathcal{F}_k$ denotes the sigma‑algebra generated by
$\{w_0,e_0,\xi_0,\dots,\xi_{k-1}\}$.
\end{lemma}
\begin{proof}[Proof of Lemma~\ref{lem:error}]
From (1) and unbiasedness of $Q$,
\[
\E\!\big[e_{k+1}\mid\mathcal{F}_k\big]=e_k+\nabla f(w_k)-\E[Q(u_k)\mid\mathcal{F}_k]=0 .
\]
Hence
\[
\E\!\big[\norm{e_{k+1}}^2\mid\mathcal{F}_k\big]
 = \E\!\big[\norm{Q(u_k)-u_k}^2\mid\mathcal{F}_k\big]
 \stackrel{\text{Ass.}\ref{ass:quant}}{\le}
 (1-\delta)\norm{u_k}^2 .
\]
Decompose $u_k=g_k+e_k$ and use $\E[g_k\mid\mathcal{F}_k]=\nabla f(w_k)$ and the
variance bound $\E\!\big[\norm{g_k-\nabla f(w_k)}^2\mid\mathcal{F}_k\big]\le\sigma^2$,
which yields the claimed inequality. \hfill$\square$
\end{proof}

\newpage

\section*{Main Proof (Step‑by‑Step Derivation)}
We prove Theorem~\ref{thm:main}.  
All expectations $\E[\cdot]$ are taken over the randomness of the
stochastic gradients $\{\xi_k\}$ and the quantizer $Q$.

\noindent\textbf{Step 1. Smoothness descent.}  
From Lemma~\ref{lem:smooth} with $x=w_k$, $y=w_{k+1}=w_k-\eta\tilde g_k$,
\[
f(w_{k+1})\le f(w_k)-\eta\ip{\nabla f(w_k)}{\tilde g_k}
               +\frac{L\eta^{2}}{2}\norm{\tilde g_k}^{2}. \tag{4}
\]

\noindent\textbf{Step 2. Expand the inner product.}  
Write $\tilde g_k = g_k+e_k-(g_k+e_k-Q(g_k+e_k)) = g_k+e_k-e_{k+1}$,
so that
\[
\ip{\nabla f(w_k)}{\tilde g_k}
 = \norm{\nabla f(w_k)}^{2}
   +\ip{\nabla f(w_k)}{e_k-e_{k+1}} . \tag{5}
\]

\noindent\textbf{Step 3. Take conditional expectation.}  
Condition on $\mathcal{F}_k$ and use unbiasedness of $g_k$ and $Q$:
\[
\E\!\big[\ip{\nabla f(w_k)}{e_k-e_{k+1}}\mid\mathcal{F}_k\big]
= \ip{\nabla f(w_k)}{e_k-\E[e_{k+1}\mid\mathcal{F}_k]}=0 . \tag{6}
\]
Hence
\[
\E\!\big[\ip{\nabla f(w_k)}{\tilde g_k}\mid\mathcal{F}_k\big]
 = \norm{\nabla f(w_k)}^{2}. \tag{7}
\]

\noindent\textbf{Step 4. Bound the second‑moment term.}  
Using unbiasedness and variance bound,
\begin{align*}
\E\!\big[\norm{\tilde g_k}^{2}\mid\mathcal{F}_k\big]
&= \E\!\big[\norm{Q(u_k)}^{2}\mid\mathcal{F}_k\big] \\
&= \E\!\big[\norm{u_k}^{2}\mid\mathcal{F}_k\big]
   -\E\!\big[\norm{u_k-Q(u_k)}^{2}\mid\mathcal{F}_k\big] \\
&\le \E\!\big[\norm{u_k}^{2}\mid\mathcal{F}_k\big]
   -(1-\delta)\E\!\big[\norm{u_k}^{2}\mid\mathcal{F}_k\big] \quad\text{[by Ass.~\ref{ass:quant}]}\\
&= \delta \E\!\big[\norm{u_k}^{2}\mid\mathcal{F}_k\big] . \tag{8}
\end{align*}
Now $u_k=g_k+e_k$, therefore
\[
\E\!\big[\norm{u_k}^{2}\mid\mathcal{F}_k\big]
 = \norm{e_k}^{2} + \E\!\big[\norm{g_k}^{2}\mid\mathcal{F}_k\big]
   +2\ip{e_k}{\nabla f(w_k)} .
\]
Using Cauchy–Schwarz and $\E\!\big[\norm{g_k}^{2}\big]\le G^{2}$ (Assumption~\ref{ass:bound}),
\[
\E\!\big[\norm{u_k}^{2}\mid\mathcal{F}_k\big]
 \le \norm{e_k}^{2}+G^{2}+2\norm{e_k}\norm{\nabla f(w_k)} . \tag{9}
\]

\noindent\textbf{Step 5. Combine (4)–(9).}  
Take total expectation of (4) and substitute (7) and (8):
\begin{align*}
\E[f(w_{k+1})]
&\le \E[f(w_k)] -\eta \E\!\big[\norm{\nabla f(w_k)}^{2}\big]
   +\frac{L\eta^{2}}{2}\E\!\big[\norm{\tilde g_k}^{2}\big] \\
&\le \E[f(w_k)] -\eta \E\!\big[\norm{\nabla f(w_k)}^{2}\big]
   +\frac{L\eta^{2}\delta}{2}\E\!\big[\norm{u_k}^{2}\big] . \tag{10}
\end{align*}
Insert bound (9) into (10):
\[
\E[f(w_{k+1})]
\le \E[f(w_k)] -\eta \E\!\big[\norm{\nabla f(w_k)}^{2}\big]
   +\frac{L\eta^{2}\delta}{2}\Big(
        \E[\norm{e_k}^{2}] + G^{2}
        +2\E[\norm{e_k}\norm{\nabla f(w_k)}]\Big). \tag{11}
\]

\noindent\textbf{Step 6. Handle the cross term.}  
By Young’s inequality $2ab\le \frac{\eta}{2}a^{2}+\frac{2}{\eta}b^{2}$,
\[
2\E[\norm{e_k}\norm{\nabla f(w_k)}]
\le \frac{\eta}{2}\E[\norm{\nabla f(w_k)}^{2}]
   +\frac{2}{\eta}\E[\norm{e_k}^{2}]. \tag{12}
\]

\noindent\textbf{Step 7. Gather the $\norm{\nabla f(w_k)}^{2}$ terms.}
Plug (12) into (11):
\begin{align*}
\E[f(w_{k+1})]
&\le \E[f(w_k)] -
      \Big(\eta-\frac{L\eta^{2}\delta}{4}\Big)
      \E[\norm{\nabla f(w_k)}^{2}] \\
&\quad +\frac{L\eta^{2}\delta}{2}G^{2}
   +\Big(\frac{L\eta^{2}\delta}{2}\cdot\frac{2}{\eta}\Big)\E[\norm{e_k}^{2}]
   +\frac{L\eta^{2}\delta}{2}\E[\norm{e_k}^{2}] .
\end{align*}
Simplify the coefficients of $\E[\norm{e_k}^{2}]$:
\[
\frac{L\eta^{2}\delta}{2}\Big(\frac{2}{\eta}+1\Big)
 = L\eta\delta +\frac{L\eta^{2}\delta}{2}. \tag{13}
\]

\noindent\textbf{Step 8. Choose stepsize to make the gradient term positive.}  
We require
\[
\eta-\frac{L\eta^{2}\delta}{4}\ge\frac{\eta}{2}
\;\Longleftrightarrow\;
\eta\le\frac{2}{L\delta}. \tag{14}
\]
A more restrictive but convenient choice is
\[
\eta\le\frac{\delta}{L(2-\delta)} , \tag{15}
\]
which satisfies (14) and also guarantees that the coefficient in front of
$\E[\norm{e_k}^{2}]$ in (13) does not exceed $\frac{\eta}{2}$ (see Lemma~\ref{lem:error}).

\noindent\textbf{Step 9. Recursion for the error norm.}  
From Lemma~\ref{lem:error} and the stepsize bound (15) we have
\[
\E[\norm{e_{k+1}}^{2}]
\le (1-\delta)\Big(\E[\norm{e_k}^{2}]
      +\E[\norm{\nabla f(w_k)}^{2}]+\sigma^{2}\Big). \tag{16}
\]
Iterating (16) and using $e_0=0$ yields a uniform bound
\[
\E[\norm{e_k}^{2}]\le \frac{1-\delta}{\delta}\big(G^{2}+\sigma^{2}\big)
=:C_{e},\qquad\forall k. \tag{17}
\]

\noindent\textbf{Step 10. Plug the uniform error bound.}  
Using (17) in the descent inequality we obtain
\[
\E[f(w_{k+1})]
\le \E[f(w_k)] -\frac{\eta}{2}\E[\norm{\nabla f(w_k)}^{2}]
   +\frac{L\eta^{2}\delta}{2}G^{2}
   +\Big(L\eta\delta+\frac{L\eta^{2}\delta}{2}\Big)C_{e}. \tag{18}
\]

\noindent\textbf{Step 11. Sum over $k=0,\dots,T-1$.}  
Telescoping the first term gives
\[
\sum_{k=0}^{T-1}\E[\norm{\nabla f(w_k)}^{2}]
\le \frac{2\big(f(w_0)-f(w^{*})\big)}{\eta}
   + L\eta\delta T G^{2}
   + 2\big(L\eta\delta+\tfrac{L\eta^{2}\delta}{2}\big)C_{e}T . \tag{19}
\]
Dividing by $T$ yields exactly inequality (3) of the theorem
(after absorbing the constants involving $C_{e}$ into the term
$\frac{L\eta G^{2}}{\delta}$, which dominates for the stepsize range
considered). \hfill$\square$

\section*{Rate Derivation (With Constants)}
From (3) set $\eta = \frac{c}{\sqrt{T}}$ with
$c\le \frac{\delta}{L(2-\delta)}$. Then
\[
\frac{2\big(f(w_0)-f(w^{*})\big)}{\eta T}
   =\frac{2\big(f(w_0)-f(w^{*})\big)}{c}\,\frac{1}{\sqrt{T}},
\qquad
\frac{L\eta G^{2}}{\delta}
   =\frac{L G^{2}c}{\delta}\,\frac{1}{\sqrt{T}} .
\]
Thus
\[
\frac{1}{T}\sum_{k=0}^{T-1}\E[\norm{\nabla f(w_k)}^{2}]
\le
\underbrace{\frac{2\big(f(w_0)-f(w^{*})\big)}{c}
           +\frac{L G^{2}c}{\delta}}_{=:C}\,
\frac{1}{\sqrt{T}} .
\]
Choosing $c=\sqrt{\frac{2\delta\big(f(w_0)-f(w^{*})\big)}{L G^{2}}}$ minimises $C$,
which leads to the explicit rate
\[
\frac{1}{T}\sum_{k=0}^{T-1}\E[\norm{\nabla f(w_k)}^{2}]
\le
2\sqrt{\frac{L G^{2}}{\delta}\big(f(w_0)-f(w^{*})\big)}\,
\frac{1}{\sqrt{T}} .
\]

\section*{Special Cases}
\begin{enumerate}
\item \textbf{Exact communication.} $\delta=1$ (i.e., $Q(x)=x$). The bound
reduces to the classical SGD rate
$\mathcal{O}\big(1/\sqrt{T}\big)$.
\item \textbf{Very coarse quantizer.} $\delta\ll1$.
The stepsize must be reduced according to (15), and the constant
prefactor grows like $1/\delta$, reflecting the loss of precision.
\item \textbf{Deterministic gradients.} $\sigma=0$.
The term $C_{e}$ in (17) simplifies to $\frac{1-\delta}{\delta}G^{2}$,
so the same $O(1/\sqrt{T})$ rate holds without stochastic variance.
\end{enumerate}

\end{document}